\chapter{Osteogenesis Imperfecta}
\index{Osteogenesis Imperfecta}

\medskip
\noindent\textit{This chapter is for general education. A fracture or breathing problem needs prompt medical assessment.}

\section*{What it is}
Osteogenesis imperfecta (OI), sometimes called brittle bone disease, is a group of inherited disorders affecting collagen and bone strength. Severity varies widely. Some people have frequent fractures from minor trauma; others have few fractures and are diagnosed later. Some forms also affect teeth, hearing, joints, growth, or the lungs.

\section*{Features and diagnosis}
Features may include repeated fractures, bone deformity, short stature, blue or grey sclerae, loose joints, dentinogenesis imperfecta, hearing loss, or family history. Diagnosis uses the clinical pattern, family history, imaging, bone density assessment when appropriate, and sometimes genetic testing. A negative genetic test does not exclude every form of OI.

\section*{Treatment and daily care}
There is no single cure. Care is coordinated by specialists and may include safe physical activity, physiotherapy, fracture treatment, orthopaedic procedures, dental and hearing care, pain management, and medicines that improve bone strength in selected people. An occupational therapist can help with transfers, equipment, and safer independence. Exercise and handling plans should be individualised; avoiding all movement can weaken bones and muscles.

\section*{When to Seek Medical Care}
Seek urgent care after a suspected fracture, sudden severe pain, new loss of movement, breathing difficulty, or head or neck injury. Discuss new hearing changes, dental problems, worsening pain, or repeated falls with the care team. A child with unexplained fractures needs a careful medical assessment that considers both OI and other causes.

\section*{Family planning and outlook}
Genetic counselling can explain inheritance, testing options, and reproductive choices. Prognosis depends on the type of OI, fracture burden, mobility, and associated problems. A personalised emergency and fracture plan can help families and clinicians respond safely.

\section*{Sources and further reading}
\begin{itemize}
\item MedlinePlus Genetics. \textit{Osteogenesis imperfecta}. \url{https://medlineplus.gov/genetics/condition/osteogenesis-imperfecta/}
\item National Institute of Arthritis and Musculoskeletal and Skin Diseases. \textit{Osteogenesis Imperfecta}. \url{https://www.niams.nih.gov/health-topics/osteogenesis-imperfecta}
\end{itemize}
